\documentclass[11pt,a4paper]{article}
\usepackage[utf8]{inputenc}
\usepackage[T1]{fontenc}
\usepackage[catalan]{babel}
\usepackage{lmodern}
\usepackage[sfdefault,lining]{FiraSans}
\usepackage[a4paper,top=2.3cm,bottom=2.2cm,left=1.8cm,right=1.8cm,headheight=1.3cm,headsep=0.4cm,footskip=1cm]{geometry}
\usepackage{amsmath,amssymb}
\usepackage{icomma}
\usepackage{graphicx}
\usepackage[table]{xcolor}
\usepackage{tikz}
\usetikzlibrary{calc}
\usepackage{fontawesome5}
\usepackage{ccicons}
\usepackage{fancyhdr}
\usepackage{lastpage}
\usepackage{tasks}
\usepackage{needspace}
\usepackage[most]{tcolorbox}
\usepackage{array}
\usepackage[hidelinks,pdftitle={Fitxa d'exercicis de la unitat 3: Inequacions i sistemes d'inequacions. 1r de Batxillerat CCSS},pdfsubject={Matemàtiques Aplicades a les CCSS I},pdfkeywords={inequacions, primer grau, segon grau, sistemes d'inequacions, semiplans, regions del pla, programació lineal}]{hyperref}

\definecolor{navy}{HTML}{1B2748}
\definecolor{accent}{HTML}{E8603A}
\definecolor{teal}{HTML}{1F9E90}
\definecolor{soft}{HTML}{EEF1F8}
\definecolor{muted}{HTML}{5D6685}

\setlength{\parindent}{0pt}
\setlength{\parskip}{3pt}
\settasks{label=\alph*), label-format=\bfseries\color{navy}, label-width=1.3em, item-indent=1.8em, column-sep=0.6em, after-item-skip=5pt, before-skip=3pt, after-skip=4pt}

% ---------- capçalera i peu ----------
\newcommand{\llicencia}{\href{https://creativecommons.org/licenses/by-nc-sa/4.0/deed.ca}{\ccbyncsa}\ \ Llicència Creative Commons Reconeixement-NoComercial-CompartirIgual 4.0 Internacional (CC BY-NC-SA 4.0)}
\pagestyle{fancy}
\fancyhf{}
\fancyhead[L]{\raisebox{-0.35\height}{\includegraphics[height=0.95cm]{logo-gva.png}}}
\fancyhead[C]{\small\color{muted}Fitxa d'exercicis · Unitat 3: Inequacions i sistemes d'inequacions · 1r Batxillerat CCSS}
\fancyhead[R]{\raisebox{-0.35\height}{\includegraphics[height=0.95cm]{logo-vermellar.png}}}
\fancyfoot[L]{\scriptsize\color{muted}\llicencia}
\fancyfoot[R]{\small\color{muted}\thepage}
\renewcommand{\headrulewidth}{0.4pt}
\renewcommand{\headrule}{\color{muted!40}\hrule height \headrulewidth}
\fancypagestyle{primera}{\fancyhf{}\renewcommand{\headrulewidth}{0pt}
  \fancyfoot[L]{\scriptsize\color{muted}\llicencia}
  \fancyfoot[R]{\small\color{muted}\thepage}}

% ---------- nivells i exercicis ----------
\newcommand{\estrelles}[1]{{\color{accent}\ifcase#1\or\faStar\faStar[regular]\faStar[regular]\or\faStar\faStar\faStar[regular]\or\faStar\faStar\faStar\fi}}
\newcounter{exer}
\newcommand{\ex}[1]{\par\needspace{4\baselineskip}\medskip\refstepcounter{exer}\noindent{\color{navy}\bfseries\theexer.}\ \ }

\newcommand{\seccio}[2]{\par\needspace{16\baselineskip}\bigskip
  \begin{tcolorbox}[colback=navy,colframe=navy,coltext=white,arc=2mm,boxrule=0pt,left=3mm,right=3mm,top=1.2mm,bottom=1.2mm]
  {\large\bfseries #1\quad #2}
  \end{tcolorbox}\vspace{-2pt}}

\newcommand{\exf}[1]{\par\medskip\refstepcounter{exer}\noindent{\color{navy}\bfseries\theexer.}\ \ }
\providecommand{\lt}{<}\providecommand{\gt}{>}
\newcolumntype{C}[1]{>{\centering\arraybackslash}p{#1}}
% quadre de recordatori (només a la fitxa d'exercicis)
\newcommand{\recorda}[1]{\par\smallskip\begin{tcolorbox}[colback=accent!7,colframe=accent!55,arc=2mm,boxrule=0.5pt,left=2.5mm,right=2.5mm,top=1mm,bottom=1mm]\small{\color{accent}\faLightbulb[regular]}\ \textbf{Recorda:}\ #1\end{tcolorbox}}
\newcommand{\res}[1]{#1}\newcommand{\rest}[1]{#1}

\def\fmneed{11}% que el títol d'una secció no quede sol al final del full
% =====================================================================
% Capçalera FentMates de les fitxes: un color per a cada curs.
% Els colors són els de la pàgina inicial (LEVELS a assets/js/materials.js);
% el logo ve de _eines/logo/pdf/ (python3 _eines/logo/genera.py).
%
% Ús, al final del preàmbul de la unitat:
%   % =====================================================================
% Capçalera FentMates de les fitxes: un color per a cada curs.
% Els colors són els de la pàgina inicial (LEVELS a assets/js/materials.js);
% el logo ve de _eines/logo/pdf/ (python3 _eines/logo/genera.py).
%
% Ús, al final del preàmbul de la unitat:
%   % =====================================================================
% Capçalera FentMates de les fitxes: un color per a cada curs.
% Els colors són els de la pàgina inicial (LEVELS a assets/js/materials.js);
% el logo ve de _eines/logo/pdf/ (python3 _eines/logo/genera.py).
%
% Ús, al final del preàmbul de la unitat:
%   \input{../comu/fentmates.tex}
%   \fmcurs{1eso}
% i, al començament del document:
%   \capcalera{Matemàtiques}{1r ESO}{Fitxa d'exercicis · Unitat 2}{Divisibilitat}
% Opcional abans de \input: \def\fmneed{14} (espai mínim abans de cada secció).
% =====================================================================
\makeatletter
\def\fm@logodir{../../logo/pdf}
\providecommand{\fmneed}{10}
\def\fm@cursnom{}

% clau del curs, color i color del text damunt del color
\newcommand{\fm@curs}[3]{\@namedef{fm@color@#1}{#2}\@namedef{fm@text@#1}{#3}}
\fm@curs{1eso}{E6B84A}{navy}
\fm@curs{2eso}{3B82F6}{white}
\fm@curs{3eso}{FF7A50}{white}
\fm@curs{4eso-a}{A78BFA}{white}
\fm@curs{4eso-b}{6D28D9}{white}
\fm@curs{1batx-mi}{9B1C1C}{white}
\fm@curs{1batx-ccss}{EF6B6B}{white}
\fm@curs{2batx-mii}{1E6B47}{white}
\fm@curs{2batx-ccss}{2BB3A3}{white}

\newcommand{\fmcurs}[1]{%
  \@ifundefined{fm@color@#1}{\PackageError{fentmates}{Curs desconegut: #1}{}}{}%
  \edef\fm@tmp{\noexpand\definecolor{curs}{HTML}{\@nameuse{fm@color@#1}}%
               \noexpand\colorlet{curstext}{\@nameuse{fm@text@#1}}}\fm@tmp
  \colorlet{cursfosc}{curs!72!black}%
  \def\fm@logo{\fm@logodir/vertical-#1.pdf}%
  \def\fm@marca{\fm@logodir/marca-#1.pdf}}

% ---------- capçalera de la primera pàgina ----------
% #1 matèria (petit), #2 curs, #3 tipus de document i unitat, #4 títol
\newcommand{\capcalera}[4]{%
  \gdef\fm@cursnom{#2}%
  \par\noindent
  \begin{tcolorbox}[colback=curs,colframe=curs,boxrule=0pt,arc=2.2mm,width=3.5cm,height=2.15cm,
      valign=center,left=3mm,right=2mm,top=1mm,bottom=1mm,nobeforeafter,box align=center]
    \color{curstext}{\footnotesize #1}\par\vspace{1pt}{\fontsize{19}{21}\selectfont\bfseries #2\par}%
  \end{tcolorbox}\hspace{5mm}%
  \begin{minipage}[c]{\dimexpr\linewidth-3.5cm-5mm-2.9cm\relax}
    {\color{cursfosc}\small\bfseries\MakeUppercase{#3}}\par\vspace{2pt}
    {\color{navy}\fontsize{20}{23}\selectfont\bfseries #4\par}
  \end{minipage}\hfill
  \begin{minipage}[c]{2.6cm}\raggedleft\includegraphics[width=2.6cm]{\fm@logo}\end{minipage}\par
  \vspace{2pt}}

% ---------- barres de secció del color del curs ----------
\renewcommand{\seccio}[2]{\par\needspace{\fmneed\baselineskip}\bigskip
  \begin{tcolorbox}[colback=curs,colframe=curs,coltext=curstext,arc=2mm,boxrule=0pt,left=3mm,right=3mm,top=1.2mm,bottom=1.2mm]
  {\large\bfseries #1\quad #2}
  \end{tcolorbox}\vspace{-2pt}}

% ---------- estrelles de dificultat del color del curs (a Batxillerat no se'n posen) ----------
\@ifundefined{estrelles}{}{%
  \renewcommand{\estrelles}[1]{{\color{curs}\ifcase#1\or\faStar\faStar[regular]\faStar[regular]\or\faStar\faStar\faStar[regular]\or\faStar\faStar\faStar\fi}}}

% ---------- capçalera i peu de les pàgines ----------
% primera pàgina: logos institucionals al peu
\fancypagestyle{primera}{\fancyhf{}\renewcommand{\headrulewidth}{0pt}
  \fancyfoot[L]{\raisebox{-0.32\height}{\includegraphics[height=0.78cm]{logo-gva.png}}\hspace{5mm}%
    \raisebox{-0.32\height}{\includegraphics[height=0.78cm]{logo-vermellar.png}}}
  \fancyfoot[R]{\scriptsize\color{muted}\href{https://creativecommons.org/licenses/by-nc-sa/4.0/deed.ca}{\ccbyncsa}\ \ CC BY-NC-SA 4.0 · fentmates.org\qquad\small\thepage}}
% la resta: el símbol a l'esquerra i el curs a la dreta, amb el títol al centre com abans
\fancyhead[L]{\raisebox{-0.3\height}{\includegraphics[height=0.62cm]{\fm@marca}}}
\fancyhead[R]{\small\bfseries\color{cursfosc}\fm@cursnom}
\renewcommand{\headrule}{\color{curs!45}\hrule height \headrulewidth}
\makeatother

%   \fmcurs{1eso}
% i, al començament del document:
%   \capcalera{Matemàtiques}{1r ESO}{Fitxa d'exercicis · Unitat 2}{Divisibilitat}
% Opcional abans de \input: \def\fmneed{14} (espai mínim abans de cada secció).
% =====================================================================
\makeatletter
\def\fm@logodir{../../logo/pdf}
\providecommand{\fmneed}{10}
\def\fm@cursnom{}

% clau del curs, color i color del text damunt del color
\newcommand{\fm@curs}[3]{\@namedef{fm@color@#1}{#2}\@namedef{fm@text@#1}{#3}}
\fm@curs{1eso}{E6B84A}{navy}
\fm@curs{2eso}{3B82F6}{white}
\fm@curs{3eso}{FF7A50}{white}
\fm@curs{4eso-a}{A78BFA}{white}
\fm@curs{4eso-b}{6D28D9}{white}
\fm@curs{1batx-mi}{9B1C1C}{white}
\fm@curs{1batx-ccss}{EF6B6B}{white}
\fm@curs{2batx-mii}{1E6B47}{white}
\fm@curs{2batx-ccss}{2BB3A3}{white}

\newcommand{\fmcurs}[1]{%
  \@ifundefined{fm@color@#1}{\PackageError{fentmates}{Curs desconegut: #1}{}}{}%
  \edef\fm@tmp{\noexpand\definecolor{curs}{HTML}{\@nameuse{fm@color@#1}}%
               \noexpand\colorlet{curstext}{\@nameuse{fm@text@#1}}}\fm@tmp
  \colorlet{cursfosc}{curs!72!black}%
  \def\fm@logo{\fm@logodir/vertical-#1.pdf}%
  \def\fm@marca{\fm@logodir/marca-#1.pdf}}

% ---------- capçalera de la primera pàgina ----------
% #1 matèria (petit), #2 curs, #3 tipus de document i unitat, #4 títol
\newcommand{\capcalera}[4]{%
  \gdef\fm@cursnom{#2}%
  \par\noindent
  \begin{tcolorbox}[colback=curs,colframe=curs,boxrule=0pt,arc=2.2mm,width=3.5cm,height=2.15cm,
      valign=center,left=3mm,right=2mm,top=1mm,bottom=1mm,nobeforeafter,box align=center]
    \color{curstext}{\footnotesize #1}\par\vspace{1pt}{\fontsize{19}{21}\selectfont\bfseries #2\par}%
  \end{tcolorbox}\hspace{5mm}%
  \begin{minipage}[c]{\dimexpr\linewidth-3.5cm-5mm-2.9cm\relax}
    {\color{cursfosc}\small\bfseries\MakeUppercase{#3}}\par\vspace{2pt}
    {\color{navy}\fontsize{20}{23}\selectfont\bfseries #4\par}
  \end{minipage}\hfill
  \begin{minipage}[c]{2.6cm}\raggedleft\includegraphics[width=2.6cm]{\fm@logo}\end{minipage}\par
  \vspace{2pt}}

% ---------- barres de secció del color del curs ----------
\renewcommand{\seccio}[2]{\par\needspace{\fmneed\baselineskip}\bigskip
  \begin{tcolorbox}[colback=curs,colframe=curs,coltext=curstext,arc=2mm,boxrule=0pt,left=3mm,right=3mm,top=1.2mm,bottom=1.2mm]
  {\large\bfseries #1\quad #2}
  \end{tcolorbox}\vspace{-2pt}}

% ---------- estrelles de dificultat del color del curs (a Batxillerat no se'n posen) ----------
\@ifundefined{estrelles}{}{%
  \renewcommand{\estrelles}[1]{{\color{curs}\ifcase#1\or\faStar\faStar[regular]\faStar[regular]\or\faStar\faStar\faStar[regular]\or\faStar\faStar\faStar\fi}}}

% ---------- capçalera i peu de les pàgines ----------
% primera pàgina: logos institucionals al peu
\fancypagestyle{primera}{\fancyhf{}\renewcommand{\headrulewidth}{0pt}
  \fancyfoot[L]{\raisebox{-0.32\height}{\includegraphics[height=0.78cm]{logo-gva.png}}\hspace{5mm}%
    \raisebox{-0.32\height}{\includegraphics[height=0.78cm]{logo-vermellar.png}}}
  \fancyfoot[R]{\scriptsize\color{muted}\href{https://creativecommons.org/licenses/by-nc-sa/4.0/deed.ca}{\ccbyncsa}\ \ CC BY-NC-SA 4.0 · fentmates.org\qquad\small\thepage}}
% la resta: el símbol a l'esquerra i el curs a la dreta, amb el títol al centre com abans
\fancyhead[L]{\raisebox{-0.3\height}{\includegraphics[height=0.62cm]{\fm@marca}}}
\fancyhead[R]{\small\bfseries\color{cursfosc}\fm@cursnom}
\renewcommand{\headrule}{\color{curs!45}\hrule height \headrulewidth}
\makeatother

%   \fmcurs{1eso}
% i, al començament del document:
%   \capcalera{Matemàtiques}{1r ESO}{Fitxa d'exercicis · Unitat 2}{Divisibilitat}
% Opcional abans de \input: \def\fmneed{14} (espai mínim abans de cada secció).
% =====================================================================
\makeatletter
\def\fm@logodir{../../logo/pdf}
\providecommand{\fmneed}{10}
\def\fm@cursnom{}

% clau del curs, color i color del text damunt del color
\newcommand{\fm@curs}[3]{\@namedef{fm@color@#1}{#2}\@namedef{fm@text@#1}{#3}}
\fm@curs{1eso}{E6B84A}{navy}
\fm@curs{2eso}{3B82F6}{white}
\fm@curs{3eso}{FF7A50}{white}
\fm@curs{4eso-a}{A78BFA}{white}
\fm@curs{4eso-b}{6D28D9}{white}
\fm@curs{1batx-mi}{9B1C1C}{white}
\fm@curs{1batx-ccss}{EF6B6B}{white}
\fm@curs{2batx-mii}{1E6B47}{white}
\fm@curs{2batx-ccss}{2BB3A3}{white}

\newcommand{\fmcurs}[1]{%
  \@ifundefined{fm@color@#1}{\PackageError{fentmates}{Curs desconegut: #1}{}}{}%
  \edef\fm@tmp{\noexpand\definecolor{curs}{HTML}{\@nameuse{fm@color@#1}}%
               \noexpand\colorlet{curstext}{\@nameuse{fm@text@#1}}}\fm@tmp
  \colorlet{cursfosc}{curs!72!black}%
  \def\fm@logo{\fm@logodir/vertical-#1.pdf}%
  \def\fm@marca{\fm@logodir/marca-#1.pdf}}

% ---------- capçalera de la primera pàgina ----------
% #1 matèria (petit), #2 curs, #3 tipus de document i unitat, #4 títol
\newcommand{\capcalera}[4]{%
  \gdef\fm@cursnom{#2}%
  \par\noindent
  \begin{tcolorbox}[colback=curs,colframe=curs,boxrule=0pt,arc=2.2mm,width=3.5cm,height=2.15cm,
      valign=center,left=3mm,right=2mm,top=1mm,bottom=1mm,nobeforeafter,box align=center]
    \color{curstext}{\footnotesize #1}\par\vspace{1pt}{\fontsize{19}{21}\selectfont\bfseries #2\par}%
  \end{tcolorbox}\hspace{5mm}%
  \begin{minipage}[c]{\dimexpr\linewidth-3.5cm-5mm-2.9cm\relax}
    {\color{cursfosc}\small\bfseries\MakeUppercase{#3}}\par\vspace{2pt}
    {\color{navy}\fontsize{20}{23}\selectfont\bfseries #4\par}
  \end{minipage}\hfill
  \begin{minipage}[c]{2.6cm}\raggedleft\includegraphics[width=2.6cm]{\fm@logo}\end{minipage}\par
  \vspace{2pt}}

% ---------- barres de secció del color del curs ----------
\renewcommand{\seccio}[2]{\par\needspace{\fmneed\baselineskip}\bigskip
  \begin{tcolorbox}[colback=curs,colframe=curs,coltext=curstext,arc=2mm,boxrule=0pt,left=3mm,right=3mm,top=1.2mm,bottom=1.2mm]
  {\large\bfseries #1\quad #2}
  \end{tcolorbox}\vspace{-2pt}}

% ---------- estrelles de dificultat del color del curs (a Batxillerat no se'n posen) ----------
\@ifundefined{estrelles}{}{%
  \renewcommand{\estrelles}[1]{{\color{curs}\ifcase#1\or\faStar\faStar[regular]\faStar[regular]\or\faStar\faStar\faStar[regular]\or\faStar\faStar\faStar\fi}}}

% ---------- capçalera i peu de les pàgines ----------
% primera pàgina: logos institucionals al peu
\fancypagestyle{primera}{\fancyhf{}\renewcommand{\headrulewidth}{0pt}
  \fancyfoot[L]{\raisebox{-0.32\height}{\includegraphics[height=0.78cm]{logo-gva.png}}\hspace{5mm}%
    \raisebox{-0.32\height}{\includegraphics[height=0.78cm]{logo-vermellar.png}}}
  \fancyfoot[R]{\scriptsize\color{muted}\href{https://creativecommons.org/licenses/by-nc-sa/4.0/deed.ca}{\ccbyncsa}\ \ CC BY-NC-SA 4.0 · fentmates.org\qquad\small\thepage}}
% la resta: el símbol a l'esquerra i el curs a la dreta, amb el títol al centre com abans
\fancyhead[L]{\raisebox{-0.3\height}{\includegraphics[height=0.62cm]{\fm@marca}}}
\fancyhead[R]{\small\bfseries\color{cursfosc}\fm@cursnom}
\renewcommand{\headrule}{\color{curs!45}\hrule height \headrulewidth}
\makeatother

\fmcurs{1batx-ccss}

\begin{document}
\thispagestyle{primera}

\capcalera{Matemàtiques CCSS}{1r Batx}{Fitxa d'exercicis · Unitat 3}{\raggedright Inequacions i sistemes d'inequacions}
\vspace{10pt}
Nom i cognoms: \rule{8.4cm}{0.4pt}\hfill Grup: \rule{1.6cm}{0.4pt}\hfill Data: \rule{2.2cm}{0.4pt}
\vspace{8pt}

\begin{tcolorbox}[colback=soft,colframe=soft,arc=2mm,left=3mm,right=3mm,top=1.5mm,bottom=1.5mm]
\small Resol els exercicis al quadern mostrant tots els passos. Escriu les solucions de les inequacions en forma d'interval i representa-les a la recta real; en les de dues incògnites, dibuixa la regió solució. En els problemes, defineix les incògnites i respon amb una frase.
\end{tcolorbox}

% =====================================================
\seccio{1}{Inequacions de primer grau}

\exf{1} Resol les inequacions següents i representa'n la solució:\begin{tasks}(3)
  \task $3x-7\gt5$
  \task $2x+5\le4x-1$
  \task $5-2x\ge11$
  \task $4(x-1)\lt2x+6$
  \task $3(2-x)\ge2(x+8)$
  \task $x-2(x+3)\gt4-3x$
\end{tasks}\recorda{Si multipliques o divideixes per un nombre negatiu, la desigualtat canvia de sentit.}

\ex{2} Resol les inequacions amb denominadors:\begin{tasks}(2)
  \task $\dfrac{x}{2}+1\le\dfrac{x}{3}+2$
  \task $\dfrac{x-1}{3}-\dfrac{x+2}{4}\gt1$
  \task $\dfrac{2x+1}{5}\ge\dfrac{x-3}{2}$
  \task $x-\dfrac{x+1}{2}\lt\dfrac{x-2}{3}$
\end{tasks}

\ex{2} Resol i interpreta el resultat:\begin{tasks}(2)
  \task $2(x+3)\gt2x+1$
  \task $3x-1\le3(x-2)$
\end{tasks}

\ex{2} Un gimnàs A cobra 30 € de matrícula i 25 € al mes, i un gimnàs B no cobra matrícula però costa 30 € al mes. A partir de quants mesos és més barat el gimnàs A?

\ex{2} Un costat d'un rectangle mesura 8 cm. Quant pot mesurar l'altre costat perquè el perímetre no supere els 30 cm?

\ex{2} Un venedor cobra 900 € fixos al mes més un 5\,\% de les vendes. Quant ha de vendre per a cobrar com a mínim 1\,400 €?


% =====================================================
\seccio{2}{Inequacions de segon grau}

\exf{1} Resol:\begin{tasks}(3)
  \task $x^2-3x-4\lt0$
  \task $x^2-x-6\ge0$
  \task $x^2-16\gt0$
  \task $x^2+3x\le0$
  \task $-x^2+4x-3\gt0$
  \task $2x^2-5x+2\le0$
\end{tasks}\recorda{Calcula les arrels, divideix la recta en intervals i estudia el signe en cada un (o pensa en la paràbola).}

\ex{2} Resol aquests casos especials:\begin{tasks}(2)
  \task $x^2+x+1\gt0$
  \task $x^2+2x+4\le0$
  \task $x^2-4x+4\gt0$
  \task $x^2-6x+9\le0$
\end{tasks}

\ex{2} Passa-ho tot a un membre i resol:\begin{tasks}(3)
  \task $x(x-2)\gt3$
  \task $(x+1)^2\le2x+5$
  \task $\dfrac{x^2}{2}-x\ge4$
\end{tasks}

\ex{2} El benefici d'una empresa, en milers d'euros, quan ven $x$ centenars d'unitats és $B(x)=-x^2+10x-16$. Quantes unitats ha de vendre per a tindre beneficis ($B(x)\gt0$)?

\ex{2} La suma d'un nombre i el seu quadrat és menor que 12. Quins valors pot prendre el nombre?


% =====================================================
\seccio{3}{Inequacions de grau superior}

\exf{2} Factoritza i resol amb una taula de signes:\begin{tasks}(2)
  \task $x^3-4x\ge0$
  \task $(x-1)(x+2)(x-3)\lt0$
  \task $x^3-x^2-2x\gt0$
  \task $x^3-6x^2+11x-6\ge0$
\end{tasks}

\ex{3} Compte amb les arrels dobles: un factor al quadrat no canvia de signe. Resol:\begin{tasks}(3)
  \task $x^3+x^2-x-1\le0$
  \task $x^2(x-3)\le0$
  \task $(x+2)^2(x-1)\gt0$
\end{tasks}

\ex{3} Resol la inequació biquadrada $x^4-5x^2+4\gt0$.


% =====================================================
\seccio{4}{Sistemes d'inequacions amb una incògnita}

\exf{2} Resol els sistemes d'inequacions:\begin{tasks}(3)
  \task $\left\{\begin{array}{l}x+3\gt0\\2x-1\le5\end{array}\right.$
  \task $\left\{\begin{array}{l}3x-2\lt x+4\\5-x\le2x+2\end{array}\right.$
  \task $\left\{\begin{array}{l}2x+1\gt x+4\\3x-5\lt x+1\end{array}\right.$
  \task $\left\{\begin{array}{l}\dfrac{x+1}{2}\ge x-1\\4x+3\gt x\end{array}\right.$
  \task $\left\{\begin{array}{l}x^2-4\lt0\\2x+1\ge0\end{array}\right.$
  \task $\left\{\begin{array}{l}x^2-5x+4\le0\\x^2-9\gt0\end{array}\right.$
\end{tasks}\recorda{La solució del sistema és la intersecció de les solucions: dibuixa-les en rectes una sobre l'altra.}

\ex{2} Troba els nombres naturals tals que el seu triple menys 5 és major que 10 i el seu doble més 1 és menor que 21.

\ex{2} Una família vol recórrer entre 400 i 600 km en un dia a una velocitat mitjana de 80 km/h. Quantes hores ha de conduir?


% =====================================================
\seccio{5}{Inequacions amb dues incògnites}

\exf{1} Representa la solució de les inequacions:\begin{tasks}(4)
  \task $x+y\le4$
  \task $2x-y\gt2$
  \task $y\lt3x$
  \task $x\ge-2$
\end{tasks}\recorda{Dibuixa la recta (contínua si és $\le,\ge$; discontínua si és $\lt,\gt$) i prova un punt que no hi estiga, com $O(0,0)$.}

\ex{1} Comprova si els punts $A(1,2)$, $B(3,1)$ i $C(-1,4)$ són solució de la inequació $2x+y\le5$.

\ex{2} Amb 20 € com a màxim volem comprar bolígrafs de 2 € i quaderns de 4 €. Escriu les inequacions que han de complir el nombre de bolígrafs, $x$, i de quaderns, $y$, i representa-les. Podem comprar 4 bolígrafs i 3 quaderns?


% =====================================================
\seccio{6}{Sistemes d'inequacions amb dues incògnites}

\exf{2} Resol els sistemes, troba els vèrtexs de la regió solució i digues si és fitada:\begin{tasks}(3)
  \task $\left\{\begin{array}{l}x+y\le4\\x\ge0,\ y\ge0\end{array}\right.$
  \task $\left\{\begin{array}{l}x+2y\le8\\x\le4\\x\ge0,\ y\ge0\end{array}\right.$
  \task $\left\{\begin{array}{l}y\le x\\x\le3\\y\ge0\end{array}\right.$
  \task $\left\{\begin{array}{l}x+y\ge2\\x\le4\\y\le3\end{array}\right.$
  \task $\left\{\begin{array}{l}2x+y\le6\\x+y\ge2\\x\ge0,\ y\ge0\end{array}\right.$
  \task $\left\{\begin{array}{l}x+y\ge3\\y\ge1\\x\ge0\end{array}\right.$
\end{tasks}\recorda{Cada vèrtex és el tall de dues rectes frontera: resol el sistema format per les seues equacions.}

\ex{2} Té solució el sistema $\left\{\begin{array}{l}x+2y\le2\\ x+2y\ge6\end{array}\right.$? Raona-ho.

\ex{3} Escriu un sistema d'inequacions la solució del qual siga el rectangle de vèrtexs $(1,1)$, $(4,1)$, $(4,3)$ i $(1,3)$, amb les vores incloses.


% =====================================================
\seccio{7}{Programació lineal}

\exf{1} Calcula el màxim de $F(x,y)=3x+2y$ en la regió definida per $x+y\le5$, $x\le3$, $x\ge0$, $y\ge0$.

\ex{2} Calcula el mínim de $F(x,y)=2x+3y$ amb les restriccions $x+y\ge4$, $x\le5$, $y\le5$, $x\ge0$, $y\ge0$.

\ex{3} Maximitza $F(x,y)=x+2y$ amb $x+2y\le8$, $x\le4$, $x\ge0$, $y\ge0$. Què hi observes?

\ex{2} Una pastisseria fa pastissos de xocolata i de nata. Cada dia en pot fer com a màxim 20, i com a màxim 12 de xocolata. Cada pastís de xocolata dona 6 € de benefici i cada pastís de nata, 4 €. Quants n'ha de fer de cada tipus per a obtindre el benefici màxim?

\ex{2} Un taller fa taules i cadires. Cada taula necessita 2 hores de fusteria i 1 de pintura, i cada cadira, 1 hora de fusteria i 1 de pintura. Cada setmana hi ha 40 hores de fusteria i 30 de pintura. El benefici és de 50 € per taula i de 30 € per cadira. Quantes taules i cadires ha de fer per a guanyar el màxim?

\ex{2} Un dietista combina dos aliments, A i B. Cada ració de A té 1 unitat de ferro i 3 de calci, i cada ració de B, 2 de ferro i 1 de calci. Cal prendre com a mínim 10 unitats de ferro i 15 de calci. Si una ració de A costa 2 € i una de B, 3 €, quantes racions de cada aliment fan la dieta més barata?

\ex{3} Una persona vol invertir com a màxim 10\,000 € en dos fons. El fons A dona un 4\,\% anual i el B, un 6\,\%. En el fons B vol posar com a màxim 6\,000 € i en el A, com a mínim 2\,000 €. Com ha de repartir els diners per a obtindre el màxim benefici? (Treballa en milers d'euros.)

\ex{3} Un institut organitza una excursió per a 400 alumnes. L'empresa de transport té 8 autobusos de 40 places i 10 de 50 places, però només 9 conductors. Llogar un autobús menut costa 600 € i un de gran, 800 €. Quants autobusos de cada tipus cal llogar perquè l'excursió siga el més barata possible?

\vfill
\begin{center}\color{muted}\small\faCheckCircle[regular]\ \ Ànim! Comprova els teus resultats amb les eines interactives dels apunts de la unitat.\end{center}

\end{document}
